

\section{Introduction}

While a rock, paper, scissors robot might seem benign in its use case at first,
it has the potential to show us some interesting aspects of human behavior and
our body’s internal perception of randomness. In ``A Meta-Analysis in Human
Behavioral Research'', Summer Armstrong notes that ``An overwhelming number of
experiments indicate that humans fail to produce random responses even when
instructed and motivated to do so.'' Our group, aiming to exploit and explore
this behavior to train a model that could eventually predict the behavior of the
human opponent by learning their perception of `random'. We of course, want to
know if this behavior is exploitable at all by a model, and if it can
effectively be done to win at rock, paper, scissors. \cite{armstrong_meta-analysis_2004}

The personal motivations for this project lie partly in the desire to create
show human random bias, but also to create an application of machine learning
that was fun and more approachable than a standard data analysis project.

The challenges in this project are numerable. To start, no data on rock, paper,
scissors could be found, requiring us to create a data collector. We also needed
to not just analyze the data, but implement a method for our model to apply that
data back into the game. We would also have to create the game.

% to implement: results summary


\section{Data Mining Task}


The data mining task for this project involved understanding and predicting
human behavior in the game of rock-paper-scissors. Since no existing datasets
were available, we had to design and implement a custom data collection system
to gather input from human players. This data served as the foundation for
analyzing patterns in decision-making and training a predictive model. The
ultimate goal was to determine whether human choices exhibit exploitable biases
and to create a model capable of leveraging these biases to improve the
computer's performance in the game.

We began by implementing a basic Aprori-based algorithm to find frequent item
sets. From there, we were able to use the data to figure out the most frequent
item. After evaluation, we determined that its performance was not great, so we
implemented a new algorithm with more sophisticated methods,



The amount of games the user will be prompted to complete can be changed within
the program, but in general, they will be required to play enough that
sufficient data can be gathered. Once this data is collected, it will be fed
into the Apriori algorithm to determine sets of frequent items. The data will
also be passed into a simple algorithm which merely determines the most frequent
item and its success rate.

Using this data, we adjust the odds of certain options being picked, by
adjusting the values in the array of odds that we stored earlier, to reflect the
outcome of these algorithms. In general, we want to adjust down the likelihood
of the computer picking a bad option, and adjust up the odds of the computer
picking a good option.

This does present a slight challenge, in that the odds are out of 1. Therefore,
adjusting one option means the others must be adjusted. Also since Apriori sets
of frequent items, it will potentially leave us with many things to adjust.

Aside from these minor challenges, the main challenges of this project will
include implementing these algorithms, and effectively implementing a persistent
form of data collection.

The main data mining question at hand is can we effectively use previous
gameplay to predict the choices of a human player.


\section{New Algorithm}

The general architecture of this algorithm is a probabilistic choice system. For
each move, it employs a number of "experts" --- subroutines that attempt to
predict what the opponent will do in a variety of ways. The main algorithm takes
all that information, weighted over time, and uses it to pick the best move.

The various "experts" are below:

\begin{itemize}
    \item "freq" --- looks at the overall tendency that the opponent has played
    each of rock, paper, scissors (i.e. just a global count, which decays over time)
    \item "markov1" --- uses Markov chains of length 1 to store a list of
    connected moves. For example, if the opponent has played rock then scissors
    before, him playing rock now will cause markov1 to predict scissors.
    \item "markov2" --- Uses Markov chains of length 2, very similar to markov1.
    \item "outcome" --- uses the outcome of the game to predict what will be
    played. For instance, if the opponent frequently plays scissors after
    winning, it would predict scissors after winning games.
    \item "copy" --- predicts the opponent will copy your last move
    \item "anticopy" --- predicts the opponent will play the counter to your last move
\end{itemize}

For every game, the main algorithm calls each of the experts, gets their
prediction, then weights it by time (so that newer games are weighted more
heavily than older ones). Once it has the weighted prediction (i.e. a
probability distribution of the three options), the algorithm randomly,
according to the distribution, plays one of the three options.



\section{Technical Approach}


% - Describe all the details of your algorithmic approach to solving this data
%   mining task and/or answering the data mining questions.
% - How are you addressing the challenges mentioned above
% - An algorithmic pseudo-code and/or a figure (block diagram) to illustrate the
%   approach will be good.

In a technical sense, our program is composed of two different approaches at solving this
same problem. That being a simple ``Most Popular Item'' method, and an ``Aggregate Expert''
approach. More detail of which will be disscussed further in this article.

When the program is started, we first have to load the previous sessions data to
keep a persistent memory. This accomplished by using Python pickles. Pickles are
a simple way to store binary data between sessions. We use it to store a list
containing each game, which will be discussed more later.

Next, the user is prompted to play a game or rock, paper, scissors. The computer
is set initially to pick completely randomly. This is done by storing the
likelihood of each option being picked in an array with 3 values which encode
the odds of a specific option being picked.

Right away, we want to create a list the win/loss ratio, this is also where we
create our list of game data that we load from a pickle as mentioned previously.
The win/loss ratio is our main data metric to track how good the program is
doing against the human player. The idea is that the higher the win/loss ratio,
the better it is at predicting the human's behavior.

We also want to define the array, from which the computer gets its random odds from,
as mentioned earlier.

Then, we start our main game loop. The program runs on a system of rounds and
games, where there is a certain number of ``games'' played in a certain number
of ``rounds''. The default for these values in our program is 5 games per round,
and 200 rounds.

As the games are played their results are stored in this data structure:

% include diagram
In English, all this is a list of a datatype we call a ``GameType'', which
is a \(3\times 2\) numpy array, which stores binary values. Each player is assigned a row,
which contains 3 of these binary values. These values encode the choice of that
player in that round. Essentially, one-hot encoding the choice. The GameType
wraps these two choices in the numpy array, and these numpy arrays are stored in
what is called a gameSeries.

The gameSeries represents each round of play, and is the level of data where
win/loss is calculated, and also the level on which the ``Most Popular Item''
approach is implemented.

The way this algorithm works, is it loops through all the games in the
game series, and uses a hash function to match those games with a pre-determined
game type (since there's only so many types of possible games in rock, paper, scissors).
From this, we can determine some information about the game, including, most important
to us in this method, whether the most frequent game in the last round was a ``winning'',
or a ``losing'' game.

If the game is determined to be a ``winning'' game, we increase the odds of the computer
able to roll that move, by a factor of the ``learning rate'', which doesn't function
like a typical learning rate, and instead, is simply the number of percentage points we
add or subtract from the likelyhood that number gets picked. The other two numbers in the
array are decreased or increased, in equal ammounts, opposite of the direction the most
frequent item went.

After all the games are played, and the win/loss ratio to view the data.
We also need to remember here to store our data back into the Python pickle to
maintain persistence.


\section{Evaluation Methodology}

The primary method for evaluating our approach is simply the win/loss rate. When
playing against a truly random opponent, one would expect the win/loss rate to
be close to 50\%; when playing against a human or other non-random opponent, the
win/loss rate should hopefully be much higher (depending on what variables the
opponent is taking into account).



\section{Results and Discussion}
As mentioned before, we had originally planned to implement the Apriori
algorithm to be able to get multiple most frequent itemsets. However, this
proved to be not only very difficult to implement, due to the shape of our data,
but also unnecessary.

As the project went on, it became increasingly difficult to justify the time taken
to get the algorithm to work, due to its limited percieved success compared to a simpler
system that only retrieved the most frequent item in a set. Apriori gets the most frequent
combinations of itemsets, but this information seemed irrelevant, so we decided to abandon
this method in favor of our simpler approach.

The ``Most Popular Item'', going up against another computer who chose random every time,
saw limited success.
